% Type of your thesis:
%	"bc" for Bachelor's
%	"mgr" for Master's
%	"phd" for PhD
%	"rig" for rigorosum
\def\ThesisType{bc}

% Language of your study programme:
%	"cs" for Czech
%	"en" for English
\def\StudyLanguage{cs}

% Thesis title in English (exactly as in the official assignment)
\def\ThesisTitle{Helper for solving Sudoku variants}

% Author of the thesis (you)
\def\ThesisAuthor{Šimon Brecher}

% Year when the thesis is submitted
\def\YearSubmitted{2026}

% Name of the department or institute, where the work was officially assigned
% (according to the Organizational Structure of MFF UK in English,
% see https://www.mff.cuni.cz/en/faculty/organizational-structure,
% or a full name of a department outside MFF)
\def\Department{Department of Theoretical Computer Science and Mathematical Logic}

% Is it a department (katedra), or an institute (ústav)?
\def\DeptType{Department}

% Thesis supervisor: name, surname and titles
\def\Supervisor{doc. Mgr. Petr Gregor, Ph.D.}

% Supervisor's department (again according to Organizational structure of MFF)
\def\SupervisorsDepartment{Department of Theoretical Computer Science and Mathematical Logic}

% Study programme (does not apply to rigorosum theses)
\def\StudyProgramme{B0613A140006}

% An optional dedication: you can thank whomever you wish (your supervisor,
% consultant, who provided you with tea and pizza, etc.)
\def\Dedication{%
\xxx{Dedication.}
}

% Abstract (recommended length around 80-200 words; this is not a copy of your thesis assignment!)
\def\Abstract{
Finding a solution for a Sudoku puzzle is relatively easy.
Unfortunately, people solve Sudoku puzzles differently from machines.
Because of that, if we want to explain how to solve a Sudoku puzzle, it is not enough to only know what the correct solution is.
The goal of this bachelor thesis is to show how to generate an explanation of the progress of solving a Sudoku puzzle.
}

% 3 to 5 keywords (recommended) separated by \sep
% Keywords are useful for indexing and searching for the theses by topic.
\def\ThesisKeywords{%
Sudoku\sep solving Sudoku \sep hint for solving Sudoku \sep Sudoku variant
}

% If any of your metadata strings contains TeX macros, you need to provide
% a plain-text version for use in XMP metadata embedded in the output PDF file.
% If you are not sure, check the generated thesis.xmpdata file.
\def\ThesisAuthorXMP{\ThesisAuthor}
\def\ThesisTitleXMP{\ThesisTitle}
\def\ThesisKeywordsXMP{\ThesisKeywords}
\def\AbstractXMP{\Abstract}

% If your abstracts are long and do not fit in the infopage, you can make the
% fonts a bit smaller by this setting. (Also, you should try to compress your abstract more.)
\def\InfoPageFont{}
%\def\InfoPageFont{\small}  % uncomment to decrease font size

% If you are studing in a Czech programme, you also need to provide metadata in Czech:
% (in English programmes, this is not used anywhere)

\def\ThesisTitleCS{Rádce pro řešení variant Sudoku}
\def\DepartmentCS{Katedra teoretické informatiky a matematické logiky}
\def\DeptTypeCS{Katedra}
\def\SupervisorsDepartmentCS{Katedra teoretické informatiky a matematické logiky}
\def\StudyProgrammeCS{B0613A140006}

\def\ThesisKeywordsCS{%
Sudoku\sep řešení Sudoku \sep nápověda řešení Sudoku \sep varianta Sudoku
}

\def\AbstractCS{%
Nalezení řešení pro Sudoku je poměrně jednoduché.
Lidé ale nejsou schopni řešit Sudoku stejným způsobem jako počítače.
Proto abychom byli schopni vysvětlit jak Sudoku vyřešit, nestačí pouze znát správné řešení.
Cílem této bakalářské práce je nalézt způsob, jak získat vysvětlení postupu řešení Sudoku.
}
